% Section 5: fine-tuning. 5.1 written (E6, H2); 5.2 a stub until E7's verdict. Sources: FINDINGS s19-s23, s25;
% CLAIMS A5, F1-F8, G1-G7, conflict 3. One sentence per line; each result sentence names its s## and CLAIMS row.
\section{Fine-tuning}
\label{sec:finetuning}

\subsection{The lead after fine-tuning}
\label{sec:ft-result}

\paragraph{Hypothesis and test.}
E6 part (b) asked whether OGC's Gram-metric VQ still beats C3DGS's own colour VQ at equal rate after C3DGS's 5,000-iteration fine-tuning, on the seven gate scenes (P\_FT, BD-PSNR of \code{ogc\_gram\_ft} against \code{c3dgs\_ft}, computed as P1 with the fine-tuned rows). % Am19 c; s21
It used E5's three criteria on fresh processes, four per scene, each fine-tuning both rows, so that the fine-tuned rows have a curve over the three colour thresholds. % Am19 c; s21
The hypothesis was formed after seeing E5's results on these scenes, so this is a pre-registered test on them and not a held-out claim. % Am19 c; s21; CLAIMS group F limits

\paragraph{Result (pre-registered).}
E6 part (b) passed as pre-registered: after fine-tuning OGC's VQ led C3DGS's own at equal rate with $\bar P_{\mathrm{FT}} = 0.0425$ dB, positive on 6 of 7 scenes, and $\bar P_{\mathrm{FT}} = 4.9\,\mathrm{SE}_{\mathrm{noise}}$ against the required 2 ($2\,\mathrm{SE}_{\mathrm{noise}} = 0.0174$ dB). % s21 "The verdict"; F1
The one negative scene was drjohnson, at $-0.0046$ dB (Table~\ref{tab:t2a}). % s21 per-scene table; F1
No scene was dropped and no primary row or value was missing. % s21 "The verdict"; F1

\begin{table}[t]
\centering
\footnotesize
\setlength{\tabcolsep}{4pt}
\caption{E6 part (b) per scene (dB). P\_FT is the pre-registered primary (9 decimals, as compared). $D_{s0}$ and $D_{s1}$ are the fine-tuned pair's protocol~ii difference at $j = 0$ with seeds 0 and 1, $D_s$ their mean and $v_s$ their sum of squared deviations. P1 replicated is a report-only secondary beside E5's P1. Generated from the committed bundles.}
\label{tab:t2a}
% FINDINGS section 21 (verdict table, secondary 1); E6 part (b), kaggle/gn_e6/S1-S3/gn6/
% generated by paper/figs/t2_e6_pft.py from committed bundles; do not edit
\begin{tabular}{lrrrrrrr}
\toprule
Scene & P\_FT (dB) & $D_{s0}$ & $D_{s1}$ & $D_s$ & $v_s$ & P1 repl. & P1 E5 \\
\midrule
bonsai & 0.062187712 & +0.079084 & +0.091531 & +0.085308 & 7.75e-05 & 0.4817 & 0.4841 \\
counter & 0.105805442 & +0.112069 & +0.036106 & +0.074087 & 2.89e-03 & 0.3197 & 0.3192 \\
kitchen & 0.012493186 & +0.083968 & +0.021124 & +0.052546 & 1.97e-03 & 0.4073 & 0.4111 \\
room & 0.056026648 & +0.069505 & +0.008662 & +0.039084 & 1.85e-03 & 0.3914 & 0.3869 \\
truck & 0.017281080 & +0.032199 & +0.031499 & +0.031849 & 2.45e-07 & 0.1484 & 0.1486 \\
drjohnson & -0.004616264 & +0.023915 & -0.011068 & +0.006424 & 6.12e-04 & 0.2757 & 0.2765 \\
playroom & 0.048257314 & +0.080564 & +0.074441 & +0.077502 & 1.87e-05 & 0.4651 & 0.4644 \\
\midrule
$\bar P$ ($n = 7$) & 0.042490731 &  &  &  &  & 0.3556 &  \\
SD\_pool & 0.032555899 &  &  &  &  &  &  \\
SE\_noise = SD\_pool / $\sqrt{14}$ & 0.008700930 &  &  &  &  & 0.0030 &  \\
2 $\times$ SE\_noise & 0.017401860 &  &  &  &  &  &  \\
scenes with P > 0 & 6 (5 needed) &  &  &  &  & 7 of 7 &  \\
\bottomrule
\end{tabular}

\end{table}

\begin{table}[t]
\centering
\small
\caption{E6 part (b), leave-one-out (descriptive, not pre-registered): the three criteria over the six scenes left when the row's scene is removed. Generated from the committed bundles.}
\label{tab:t2b}
% FINDINGS section 22 a (leave-one-out, descriptive); E6 part (b), kaggle/gn_e6/S1-S3/gn6/
% generated by paper/figs/t2_e6_pft.py from committed bundles; do not edit
\begin{tabular}{lrrrll}
\toprule
Left out & $\bar P$ & SE\_noise & $\bar P$ / SE\_noise & Positive & Criteria 1 / 2 / 3 \\
\midrule
bonsai & 0.039208 & 0.010098 & 3.88 & 5 of 6 (5 needed) & yes / yes / yes \\
counter & 0.031938 & 0.007936 & 4.02 & 5 of 6 (5 needed) & yes / yes / yes \\
kitchen & 0.047490 & 0.008696 & 5.46 & 5 of 6 (5 needed) & yes / yes / yes \\
room & 0.040235 & 0.008794 & 4.58 & 5 of 6 (5 needed) & yes / yes / yes \\
truck & 0.046692 & 0.010151 & 4.60 & 5 of 6 (5 needed) & yes / yes / yes \\
drjohnson & 0.050342 & 0.009723 & 5.18 & 6 of 6 (5 needed) & yes / yes / yes \\
playroom & 0.041530 & 0.010138 & 4.10 & 5 of 6 (5 needed) & yes / yes / yes \\
\bottomrule
\end{tabular}

\end{table}

\paragraph{How firm the result is (descriptive, not pre-registered).}
Leaving out any one scene (Table~\ref{tab:t2b}), every case still met all three criteria, with $\bar P / \mathrm{SE}_{\mathrm{noise}}$ from 3.88 to 5.46, but 6 of the 7 cases met the count with exactly 5 positive scenes, so the count is the tight criterion. % s22 a; F5
Across scenes P\_FT had a standard deviation of 0.0373 dB ($t = 3.0$), which $\mathrm{SE}_{\mathrm{noise}}$ does not measure. % s22 a; F5
Under the standard error that matches one curve per scene, $\mathrm{SD}_{\mathrm{pool}}/\sqrt{n}$, $\bar P_{\mathrm{FT}}$ lay 3.45 standard errors above zero against the registered 4.9, and the leave-one-out cases 2.75 to 3.86; criterion 3 still held in every case, and the registered verdict stands under its pre-registered rule. % s25; F1, F5
With OGC's bytes set to C3DGS's at each $j$, a counterfactual and not equal rate, the mean was $+0.0458$ dB against $\bar P_{\mathrm{FT}} = +0.0425$ dB, so the result is not produced by where the bytes fall. % s22 c; F7

\paragraph{Against the lead before fine-tuning (descriptive).}
As a ratio of means P\_FT was 0.1195 of P1 replicated, with per-scene ratios from $-0.017$ to $+0.331$; the two measures average over different byte ranges, so the ratio is not the share of the earlier lead that survives at any one rate. % s22 b; F6
E5's own $j = 0$ comparison after fine-tuning, $+0.0604$ dB at unequal bytes, measures a different quantity on separate runs; neither value replaces the other. % s19 secondaries; A5; CLAIMS conflict 3

\paragraph{Where along the curve (descriptive, not pre-registered).}
On all seven gate scenes the fine-tuned rows' PSNR difference at the same $j$ rose from $j = -1$ to $j = +1$: from $-0.119$ to $+0.023$ dB at $j = -1$, $+0.024$ to $+0.112$ dB at $j = 0$, and $+0.061$ to $+0.379$ dB at $j = +1$. % s22 d; F8
These points are not rate-matched: at $j = +1$ OGC's row had 3.5\% to 7.9\% more bytes, and at $j = -1$ 0.3\% to 7.9\% fewer. % s22 d; s21 secondary 2; F8
After fine-tuning OGC's \code{.npz} was smaller than C3DGS's at $j = -1$ on every scene and larger at $j = +1$ on every scene, while before fine-tuning it was larger at every point. % s21 secondary 2; F3
The rising pattern is scoped to the gate scenes: on the development scenes it rose on flowers but fell on train and was not monotone on treehill. % s23 e; F8, G5

\paragraph{The development scenes (descriptive).}
E6 part (a) ran the same rows on flowers, treehill and train with three predictions and no criterion (Table~\ref{tab:t3}). % Am19 b; s23; CLAIMS group G limits
All three predictions held on all three scenes: P1 $> 0$, P2 $> 0$, and P\_FT below P1. % s23 "The predictions"; G1
P\_FT was $+0.0341$ dB on flowers and negative on treehill ($-0.0262$ dB) and train ($-0.0066$ dB), so the third prediction held on treehill and train because P\_FT was negative. % s23 "The predictions"; G1, G2
Treehill's P\_FT stayed negative with the $j = 0$ seed-1 process in place of seed 0 ($-0.0423$ dB). % s23 b; G2
Train's P\_FT is not interpreted: its fine-tuned points moved by up to 0.19 dB between runs, and replacing its $j = 0$ fine-tuned point with the pilot's changed it by $+0.2274$ dB. % s23 a; G3
These are development scenes; no part (a) value is pooled with the gate scenes' results. % Am19 b; CLAIMS group G limits

\begin{table}[t]
\centering
\small
\caption{E6 part (a) on the development scenes (BD-PSNR, dB). Descriptive: three predictions checked, no criterion. Generated from the committed bundles.}
\label{tab:t3}
% FINDINGS section 23, The predictions (E6 part (a), descriptive; kaggle/gn_e6/S4-S5/gn6/)
% generated by paper/figs/t3_e6_part_a.py from committed bundles; do not edit
\begin{tabular}{lrrrrrl}
\toprule
Scene & P1 & P2 & scalar vs plain & plain vs c3dgs & P\_FT & Predictions 1 / 2 / 3 \\
\midrule
flowers & +0.123913 & +0.069250 & +0.068761 & -0.013664 & +0.034106 & held / held / held \\
treehill & +0.025062 & +0.002488 & +0.102726 & -0.077859 & -0.026188 & held / held / held \\
train & +0.129723 & +0.064938 & +0.008808 & +0.057308 & -0.006562 & held / held / held \\
\bottomrule
\end{tabular}

\end{table}

\paragraph{The curves.}
The right panels of Figure~\ref{fig:f1} show the fine-tuned curves; each scene's two panels share the PSNR axis, and the shaded band marks the byte range over which P\_FT averages. % paper/figs/f1_rd_curves.py; s21

\subsection{Why the lead shrinks (E7, pending)}
\label{sec:ft-why}

\TODO{E7 (Amendment 20): the within-process dissection of the shrink (arms R, G, CG, K, O). Written after all eleven sessions and E7's verdict; no mechanism is stated before then.}
